% Beamer presentation for MediChart EMR v2
\documentclass{beamer}
\usetheme{Madrid}
\usecolortheme{seagull}
\usepackage[utf8]{inputenc}
\usepackage{graphicx}
\usepackage{hyperref}

\title{MediChart EMR v2 — Ringkasan Proyek}
\author{[Nama Kelompok]}
\institute{Konsep Teknologi Informasi}
\date{[Tanggal]}

\begin{document}

\begin{frame}
  \titlepage
\end{frame}

\begin{frame}{Outline}
  \tableofcontents[hideallsubsections]
\end{frame}

\section{Pendahuluan}
\begin{frame}{Pendahuluan}
  \begin{itemize}
    \item Tujuan: memperluas MediChart v0 dengan RBAC dan Diagnostic Reports
    \item Scope: SPA client-side, penyimpanan `localStorage`, sesi di `sessionStorage`
  \end{itemize}
\end{frame}

\section{Deskripsi Sistem}
\begin{frame}{Deskripsi Sistem}
  \begin{itemize}
    \item Peran: Dokter, Perawat, Admin
    \item Fitur utama: autentikasi, manajemen pasien, SOAP, Diagnostic Reports, admin
    \item Deliverables: kode demo, laporan (Markdown/LaTeX), PlantUML diagrams, slides
  \end{itemize}
\end{frame}

\section{Arsitektur}
\begin{frame}{Arsitektur Sistem}
  \begin{itemize}
    \item SPA monolitik di `medichart.html`
    \item Modul: Auth, Patient, Clinical, DiagReport, Admin
    \item Penyimpanan: `localStorage` (data), `sessionStorage` (sesi)
  \end{itemize}
\end{frame}

\section{Desain Data}
\begin{frame}{Desain Data — Entity Highlights}
  \begin{itemize}
    \item `User`, `Patient`, `Encounter`, `ClinicalNote`, `Problem`, `DiagnosticReport`
    \item DiagnosticReport: `reportId`, `patientId`, `encounterId`, `type`, `date`, `status`, `interpretation`
  \end{itemize}
\end{frame}

\section{Keamanan}
\begin{frame}{Keamanan (RBAC)}
  \begin{itemize}
    \item Tiga peran: Dokter (full), Perawat (terbatas), Admin (manajemen)
    \item Kontrol tampilan dan aksi pada UI berdasarkan `canAccess()`
  \end{itemize}
\end{frame}

\section{Diagnostic Reports}
\begin{frame}{Diagnostic Reports (Lab \\ \& Imaging)}
  \begin{itemize}
    \item Tipe `lab` dan `imaging`; status: `pending`, `final`, `amended`
    \item Integrasi: dapat dikaitkan ke `encounterId` atau berdiri sendiri
    \item Tampilan: tab terpisah, daftar ringkasan, detail lengkap dengan interpretasi
  \end{itemize}
\end{frame}

\section{Antarmuka}
\begin{frame}{Antarmuka dan Alur}
  \begin{itemize}
    \item Login → AppShell → modul berdasarkan peran
    \item Panel: Patients, Diag Reports (Lab/Imaging), Admin
  \end{itemize}
\end{frame}

\section{Diagram Utama}
\begin{frame}{Diagram Utama}
  \begin{itemize}
    \item Class Diagram: `docs/class-diagram.puml` (render ke PNG untuk disertakan)
    \item Package Diagram (simple/detailed): `docs/package-diagram.puml`, `docs/package-diagram-detailed.puml`
    \item Component / DFD: `docs/component-diagram.puml`, `docs/dfd-level0.puml` / `docs/dfd-level1.puml`
  \end{itemize}
  \vspace{6pt}
  \footnotesize{Catatan: render PlantUML sources ke PNG/SVG lalu letakkan di `docs/diagrams/` dan gunakan `\includegraphics` untuk menampilkan diagram di PDF.}
\end{frame}

\section{Implementasi dan Validasi}
\begin{frame}{Implementasi dan Validasi}
  \begin{itemize}
    \item Implementasi: `medichart.html` (self-contained)
    \item Validasi: smoke tests, pemeriksaan keberadaan artefak, propery checks (opsional)
  \end{itemize}
\end{frame}

\section{Kesimpulan}
\begin{frame}{Kesimpulan}
  \begin{itemize}
    \item MediChart v2 meningkatkan modelling klinis melalui RBAC dan Diagnostic Reports
    \item Dokumentasi lengkap tersedia: laporan, makalah, PlantUML diagrams, slides
  \end{itemize}
\end{frame}

\begin{frame}{Pertanyaan}
  \centering Terima kasih. Pertanyaan?
\end{frame}

\end{document}
